\documentclass[11pt]{amsart}

\usepackage{amsmath,amssymb,amsthm,enumitem,hyperref}
\usepackage[margin=1.05in]{geometry}
\usepackage[T1]{fontenc}
\usepackage{lmodern}
\usepackage{microtype}

\hypersetup{colorlinks=true, linkcolor=blue, citecolor=blue, urlcolor=blue}

\newtheorem{theorem}{Theorem}
\newtheorem{lemma}[theorem]{Lemma}
\newtheorem{proposition}[theorem]{Proposition}
\newtheorem{corollary}[theorem]{Corollary}
\theoremstyle{definition}
\newtheorem*{definition}{Definition}
\newtheorem*{question}{Question}
\newtheorem*{example}{Example}
\theoremstyle{remark}
\newtheorem{remark}{Remark}[section]

\newcommand{\C}{\mathcal C}
\newcommand{\D}{\mathcal D}
\newcommand{\B}{\mathcal B}
\newcommand{\R}{\mathcal R}
\newcommand{\F}{\mathcal F}
\newcommand{\res}{\operatorname{res}}
\newcommand{\wt}{\operatorname{wt}}
\newcommand{\mult}{\operatorname{mult}}
\newcommand{\len}{\ell}


\begin{document}


A partition is a weakly decreasing finite list of nonnegative integers.  The \emph{weight} of a partition is the sum of its parts.

\begin{definition}[The Andrews--Dhar partition families]\label{def:families}
Suppose that $n$ is a positive integer.

\smallskip
\noindent
(a) Let $\C_3(n)$ be the set of partitions $\lambda$ of $n$ such that
\begin{enumerate}[label=(\roman*)]
\item the largest part of $\lambda$ is divisible by $3$, say it is $3J$;
\item every part of $\lambda$ which is at most $J$ occurs at most twice.
\end{enumerate}
Let $C_3(n):=|\C_3(n)|$.
\end{definition}

First, let $\B_3^{(2)}(N)$ be the set of partitions of $N$ into positive parts not divisible by $3$ whose largest part is congruent to $2$ modulo $3$. In  particular, we have $\B_3^{(2)}(0)=\varnothing$.  The finite ``Glaisher map''
\begin{equation}
        \Gamma:\C_3(n)\longrightarrow \B_3^{(2)}(n-1)
\end{equation}
is defined as follows.  Given $\lambda\in\C_3(n)$, let its largest part be $3J$.  Replace one copy of this largest part by $3J-1$.  For every other original part $t$, write uniquely
\[
       t=3^a u,\qquad 3\nmid u,
\]
and replace that one part $t$ by $3^a$ copies of $u$.  After sorting, the result is $\Gamma(\lambda)$.  The special part \(3J-1\) is a largest part of the output and is congruent
to \(2\) modulo \(3\). Indeed, every other original part is at most \(3J\);
parts strictly smaller than \(3J\) become parts at most \(3J-1\), while any
remaining copies of \(3J\) are converted into parts at most \(J\).

The first step is a finite version of Glaisher carrying in which one distinguished largest part is lowered by one.  The truncation at the largest part is the only point requiring a check.

For a fixed $J\ge1$, let $\C_{3,J}(n)$ be the subset of $\C_3(n)$ consisting of partitions whose largest part is $3J$.  Let $\B_{3,J}^{(2)}(n-1)$ be the subset of $\B_3^{(2)}(n-1)$ consisting of partitions whose largest part is exactly $3J-1$.

The next proposition records that this finite carrying procedure loses no information.

\begin{proposition}\label{prop:gamma}
The finite Glaisher map restricts to a bijection
\[
        \Gamma_J:\C_{3,J}(n)\longrightarrow\B_{3,J}^{(2)}(n-1).
\]
Consequently, we have
\[
        \Gamma:\C_3(n)\longrightarrow\B_3^{(2)}(n-1)
\]
is a bijection.
\end{proposition}

\begin{proof}
The construction of $\Gamma_J$ lowers the weight by $1$, since the distinguished part $3J$ is replaced by $3J-1$, while every other replacement preserves weight:
\[
        3^a u=\underbrace{u+u+\cdots+u}_{3^a\text{ copies}}.
\]
All produced parts are not divisible by $3$, and the distinguished part $3J-1$ is the largest part and is congruent to $2$ modulo $3$.  Hence $\Gamma_J$ lands in $\B_{3,J}^{(2)}(n-1)$.

We now construct the inverse.  Start with $\rho\in\B_{3,J}^{(2)}(n-1)$.  Replace one copy of the largest part $3J-1$ by $3J$.  For every $u$ with $3\nmid u$, let $M_u$ be the number of remaining copies of $u$, and put
\[
        E_u:=\max\{e\ge0:u3^e\le3J\}.
\]
Write $M_u$ uniquely as
\[
        M_u=Q_u3^{E_u}+\sum_{e=0}^{E_u-1}d_{u,e}3^e,
        \qquad Q_u\ge0,
        \qquad d_{u,e}\in\{0,1,2\}.
\]
Insert $Q_u$ copies of $u3^{E_u}$ and, for $0\le e<E_u$, insert $d_{u,e}$ copies of $u3^e$.  Do this independently for all $u$ not divisible by $3$, and sort.

The resulting partition has largest part $3J$.  It remains to check the multiplicity condition below $J$.  By the definition of $E_u$,
\[
        u3^{E_u}\le3J<u3^{E_u+1}.
\]
If $u3^{E_u}\le J$, then multiplying by $3$ gives $u3^{E_u+1}\le3J$, a contradiction.  Hence $u3^{E_u}>J$.  This top bucket may occur with arbitrary multiplicity $Q_u$.  Every lower bucket $u3^e$ with $e<E_u$ is at most $J$, and it occurs with multiplicity $d_{u,e}\in\{0,1,2\}$.  Thus every part at most $J$ occurs at most twice.

For each fixed $u$, the forward map replaces one copy of $u3^e$ by $3^e$ copies of $u$, and the inverse recovers the multiplicities by the displayed base-$3$ expansion with the top bucket separated.  The distinguished part is reversed by $3J\leftrightarrow 3J-1$.  Therefore the two maps are inverse bijections.
\end{proof}


\end{document}
